%================================================================
%  Part I: afterglow reheating, where the case stands. Every number
%  is printed in the chapter cited next to it, or follows from such
%  numbers by the arithmetic shown. No hyperlinks inside table cells.
%================================================================

\chapter{The question and the verdict}
\label{p1:chap:verdict}

%================================================================
\section{The proposal}
\label{p1:sec:proposal}
%================================================================

{\noindent}(1) An electron beam pumps a KrF laser gas by making ion pairs and
excited atoms, and KrF$^*$ forms from them, much of it through the harpoon
reaction of excited krypton with F$_2$. When the beam stops, the electrons it
leaves behind cool and are lost to attachment and recombination. This is the
afterglow. Part~\ref{p1:part:notes} describes it in the neon-buffered fill of an
electron-beam-pumped laser, Ne/Kr/F$_2$ $=0.948/0.05/0.002$ at 2~atm, by
electrons at 0.63~eV and $10^{20}$~m$^{-3}$, against 1.5~eV and
$2.72\times10^{21}$~m$^{-3}$ with the beam on (Sec.~\ref{nl:sec:regime}). The
afterglow holds 27 times fewer electrons, and colder ones.

{\noindent}(2) Reheating sends an electromagnetic wave into the afterglow. The wave
heats the free electrons by inverse bremsstrahlung (IB), the hotter electrons
excite krypton again, and the excited krypton forms KrF$^*$ while F$_2$ lasts. The
question of this compendium is whether that can make the laser more efficient,
that is, raise the laser energy per unit of electrical energy.

{\noindent}(3) It can only if five links hold (Fig.~\ref{p1:fig:chain}):
\begin{itemize}
\item[L1] \emph{Coupling.} The wave enters the gas and is absorbed through the
  depth of the gain medium, rather than reflected or absorbed in a thin skin.
\item[L2] \emph{Absorption.} The absorption coefficient that decides L1 is known.
\item[L3] \emph{Heating.} An available intensity raises the electron temperature,
  in a regime that the kinetic equation of the electrons describes.
\item[L4] \emph{Partition.} The energy the electrons take from the wave forms
  KrF$^*$ at a useful efficiency.
\item[L5] \emph{Laser and wall plug.} Laser photons extract the KrF$^*$, and the
  source of the wave is efficient enough for the gain to pay.
\end{itemize}

{\noindent}(4) A link is called \emph{verified} here when the calculations of
Part~\ref{p1:part:notes} establish it, with their inputs tested against measurements
where the notes cite such tests. Nothing in Part~\ref{p1:part:notes} measures
reheating itself. A link or a lever is \emph{invalidated} when those calculations
rule it out, and \emph{open} when they do not settle it, in which case
Chapter~\ref{p1:chap:settle} names what would.

%================================================================
\section{The verdict}
\label{p1:sec:verdict}
%================================================================

{\noindent}(1) The heating half of the proposal stands, and the efficiency half is
open. Coupling, absorption and heating are verified in silico, with one constraint
that shapes the proposal: only the THz band heats the afterglow through
centimetres, while microwaves heat a layer millimetres deep. The partition favours
the field. At the afterglow's density a wave far above the collision frequency
forms KrF$^*$ at 1.4 to 1.8 times the electron beam's rate per absorbed joule, but
that count leaves out the electron collisions with Kr$^*$ and KrF$^*$ on which the
sign of the gain rests, and no note carries the KrF$^*$ to laser energy or to the
wall plug. Two levers are invalidated: the nonlinearity of IB, and the laser's own
field. Table~\ref{p1:tab:verdict} collects the verdict, and
Chapter~\ref{p1:chap:evidence} gives the evidence link by link.

\begin{figure}[t]
\centering
\begin{tikzpicture}[>=Stealth,
  link/.style={draw, rounded corners=2pt, align=center, text width=2.6cm,
    minimum height=2.5cm, font=\small},
  ver/.style={link, fill=green!12},
  opn/.style={link, fill=orange!20},
  inv/.style={draw, rounded corners=2pt, align=center, text width=10.5cm,
    font=\small, fill=red!10}]
\node[ver] (L1) {\textbf{L1 coupling}\\[2pt] wave absorbed\\ over the depth\\[2pt] \emph{verified}};
\node[ver, right=0.36cm of L1] (L2) {\textbf{L2 absorption}\\[2pt] coefficient\\ known\\[2pt] \emph{verified}};
\node[ver, right=0.36cm of L2] (L3) {\textbf{L3 heating}\\[2pt] field sets $T_e$\\ at reachable $I$\\[2pt] \emph{verified}};
\node[opn, right=0.36cm of L3] (L4) {\textbf{L4 partition}\\[2pt] energy forms\\ KrF$^*$\\[2pt] \emph{open}};
\node[opn, right=0.36cm of L4] (L5) {\textbf{L5 laser and}\\ \textbf{wall plug}\\[2pt] KrF$^*$ extracted,\\ source pays\\[2pt] \emph{open}};
\foreach \a/\b in {L1/L2,L2/L3,L3/L4,L4/L5} \draw[->,thick] (\a) -- (\b);
\node[inv, below=0.55cm of L3] (N) {\textbf{Invalidated as levers:} the nonlinearity of
  inverse bremsstrahlung (Langdon, Silin), and the laser's own field};
\end{tikzpicture}
\caption{The chain from the wave to the laser energy, with the standing of each link
in Part~\ref*{p1:part:notes}. The THz band carries L1. The open links hinge on the
electron collisions of Kr$^*$ and KrF$^*$, the extraction of the KrF$^*$ and the
efficiency of the source.}
\label{p1:fig:chain}
\end{figure}

\begin{table}[t]
\caption{Where afterglow reheating stands, link by link. The numbers are for the
afterglow of Chapter~\ref*{nl:chapter}, $10^{20}$~m$^{-3}$, unless stated otherwise.}
\label{p1:tab:verdict}
\centering
\small
\begin{tabular}{@{}L{1.9cm} L{1.7cm} L{8.4cm} L{3.0cm}@{}}
\toprule
link & standing & what Part~\ref*{p1:part:notes} establishes & where \\
\midrule
L1 coupling & verified & 10~cm is absorbed at 0.64 to 3.8~THz, 2.6~cm at 0.67~THz, and a
microwave within millimetres. Reflection does not bind. & Secs.~\ref*{env:sec:spectrum},
\ref*{nl:sec:regime}, Tables~\ref*{env:tab:main}, \ref*{qib:tab:fbest}\\
\addlinespace
L2 absorption & verified & The classical kernel is within 1.6\,\% of the quantum one at
0.3 to 4~THz. Krypton's cross section is tested to a few per cent above 0.95~eV and to
about 20\,\% at 0.63~eV. Ions matter only in a cold afterglow. &
Secs.~\ref*{ea:sec:thermal}, \ref*{env:sec:sources}, \ref*{ea:sec:ei}\\
\addlinespace
L3 heating & verified & $1.4\times10^{3}$ to $1.5\times10^{4}$~W\,cm$^{-2}$ hold the
electrons at 0.63~eV, 5.8~kW\,cm$^{-3}$, far below Langdon's and Silin's thresholds. &
Sec.~\ref*{nl:sec:regime}, Table~\ref*{nl:tab:regime}\\
\addlinespace
nonlinearity, laser's field & invalidated & The lasing field is $7.1\times10^{3}$ times
below the heating threshold, and Langdon's regime needs 33 to 80 times it. &
Secs.~\ref*{nl:sec:answer}, \ref*{nl:sec:regime}\\
\addlinespace
L4 partition & open & 0.274 to 0.352 KrF$^*$ per absorbed joule at high frequency against
the beam's 0.198, without the electron collisions of Kr$^*$ and KrF$^*$ or the 5s to 5p
cycle. & Tables~\ref*{nl:tab:share}, \ref*{nl:tab:source}, Secs.~\ref*{nl:sec:wallplug},
\ref*{qib:sec:ladder}\\
\addlinespace
L5 laser, wall plug & open & No note follows the photons. The source must reach 0.56 to
0.72 of the beam's efficiency times the ratio of the two extractions. &
Secs.~\ref*{qib:sec:stands}, \ref*{nl:sec:wallplug}, Eq.~(\ref*{p1:eq:wallplug})\\
\bottomrule
\end{tabular}
\end{table}

%================================================================
\section{Reading Part~II}
\label{p1:sec:reading}
%================================================================

{\noindent}(1) Part~\ref{p1:part:notes} is six notes, one chapter each, in their
reading order. Chapter~\ref{env:chapter} derives the best frequency for krypton gases
and settles THz against GHz (L1, L2). Chapter~\ref{qib:chapter} follows the excited
atoms inside the electrons' balances and marches a pulse (L1, L4).
Chapter~\ref{nl:chapter} locates the lasers and the sources against the
nonlinearities of IB and counts what field pumping buys (L3, L4, L5).
Chapter~\ref{lad:chapter} tests the absorption model against measured absorption in
laser plasmas (L2). Chapter~\ref{pd:chapter} derives the quantum theory of IB on ions
and Chapter~\ref{ea:chapter} that of absorption on atoms at THz frequencies (L2, L3).

{\noindent}(2) Each chapter keeps its note's text, notation and references, with its
own numbering inside the chapter. Chapters~\ref{env:chapter}, \ref{lad:chapter},
\ref{pd:chapter} and \ref{ea:chapter} use Gaussian units and Chapters~\ref{qib:chapter}
and \ref{nl:chapter} SI units. Electron densities are in cm$^{-3}$ in
Chapters~\ref{env:chapter}, \ref{qib:chapter} and \ref{lad:chapter} and in m$^{-3}$ in
Chapters~\ref{nl:chapter} and \ref{ea:chapter}, and $10^{14}$~cm$^{-3}$ is
$10^{20}$~m$^{-3}$. ``These notes'' in a chapter means that chapter.


\chapter{The evidence, link by link}
\label{p1:chap:evidence}

%================================================================
\section{L1: coupling}
\label{p1:sec:L1}
%================================================================

{\noindent}(1) A wave heats a depth $L$ of gas best at the frequency whose absorption
length is $L$, provided the gas reflects no more than 5\,\% of it,
Eq.~\eqref{env:eq:fbest}. Far above the collision and plasma frequencies the
absorption coefficient falls as $f^{-2}$, Eq.~\eqref{env:eq:alpha}, and the frequency
that matches the depth grows as the square root of the electron density, the
collision frequency and the depth, Eq.~\eqref{env:eq:fA}. The afterglow's fewer and
colder electrons therefore need a lower frequency than the gas with the beam on.

{\noindent}(2) At the afterglow's density, $10^{14}$~cm$^{-3}$, a depth of 10~cm calls
for 1.1, 2.2 and 3.8~THz in pure krypton at 0.1, 0.5 and 2~atm
(Table~\ref{env:tab:main}), and for 0.64 to 2.4~THz in the Ne fill and 1 to 3.5~THz in
the Ar fill from 0.1 to 2~atm (Table~\ref{qib:tab:fbest}), with the electrons heated by
300~kW\,cm$^{-3}$ in both tables. In the afterglow itself, at 0.63~eV in the Ne fill, a
0.67~THz wave is absorbed in 2.6~cm (Sec.~\ref{nl:sec:regime}). There the electrons'
collision frequency is 0.12 of the wave's angular frequency (Table~\ref{nl:tab:regime})
and the cutoff lies at 0.090~THz (Sec.~\ref{env:sec:reflection}), so the absorption
coefficient falls as $f^{-2}$, and 10~cm takes $(10/2.6)^{1/2}=2.0$ times that
frequency, 1.3~THz.

{\noindent}(3) Microwaves do not reach that depth. In the afterglow a 2.45~GHz wave
lies far below the collision frequency, $\langle\nu_m\rangle/\omega=32$
(Table~\ref{nl:tab:regime}), and a 170~GHz wave heats a layer 2.5~mm deep
(Sec.~\ref{nl:sec:regime}). In krypton at $10^{14}$~cm$^{-3}$ and 0.5~atm, where the
collisions outrun the cutoff, a GHz wave does enter the gas, 43\,\% of it reflected at
2.45~GHz and 1\,\% at 95~GHz, but it is absorbed in 1.8 to 4.2~mm, and over the
pressures and both gases in 0.6 to 8.3~mm, so that the front millimetres take all of it
(Sec.~\ref{env:sec:spectrum}).

{\noindent}(4) The cutoff does not bind in the afterglow. The plasma frequency is
0.090~THz at $10^{14}$~cm$^{-3}$ (Sec.~\ref{env:sec:reflection}), where with the beam on
it is 0.47~THz (Sec.~\ref{nl:sec:share}).

\medskip
\noindent\textbf{Standing: verified, with a band.} A wave of 0.64 to 3.8~THz, 1.3~THz in
the Ne fill's afterglow, deposits its power through 10~cm of gas at the afterglow's
density. A 0.67~THz wave deposits it through 2.6~cm, and a microwave in a layer
millimetres deep. Volumetric reheating is a THz proposal.

%================================================================
\section{L2: absorption}
\label{p1:sec:L2}
%================================================================

{\noindent}(1) The absorption coefficient is the Drude form with the collision
frequency $\nueff=\nuea+\nuT\Lambda$, Eq.~\eqref{env:eq:nueff}, the atoms entering
through the energy-weighted average of their momentum-transfer collision frequency,
Eq.~\eqref{env:eq:nuea}, and the ions through the Spitzer frequency and the Gaunt
factor. Three things decide whether it is right: the kernel of the atoms, the cross
sections that the kernel takes, and the term of the ions.

{\noindent}(2) \emph{The kernel.} The classical kernel treats the electron as a
particle that absorbs continuously. Against the quantum rate between exact scattering
states it is within 1.6\,\% in either fill from 0.3 to 4~THz at every electron
temperature from 300~K to 1.08~eV (Sec.~\ref{ea:sec:thermal}), so the absorption lengths
of L1 stand. It fails only far above that band: at 28.3~THz and 300~K the quantum
absorption is 10\,\% below it in the neon fill and 28\,\% below it in the argon fill.

{\noindent}(3) \emph{The cross sections.} Krypton's momentum-transfer cross section is
tested by the drift velocities it predicts. Between 0.95 and 2.6~eV the computed drift
velocity is 0.96 to 1.02 times Nakamura's measurements \cite{p1:Nakamura1991} and 0.84
to 0.97 times those of Pack, Voshall and Phelps \cite{p1:Pack1962}
(Sec.~\ref{env:sec:sources}). Below 0.95~eV only the latter exist, and across the
Ramsauer minimum the computed drift velocity is 0.80 to 1.20 times them. Since the
drift velocity weights the inverse of the cross section over the same electrons, the
cross section is tested there to about 20\,\%, and the afterglow's 0.63~eV electrons
sit in that minimum, which lies at 0.5~eV (Sec.~\ref{env:sec:atoms}).

{\noindent}(4) \emph{The ions.} The ions matter only in a cold afterglow. While the
electrons stay at 0.63~eV the ions catch up with the atoms only above
$10^{22}$~m$^{-3}$. At 300~K they take over at 1~THz above $7.0\times10^{19}$~m$^{-3}$
in the neon fill and $1.5\times10^{20}$~m$^{-3}$ in the argon fill with the classical
Gaunt factor, and at $10^{20}$~m$^{-3}$ they absorb 1.4 times as much as the atoms in
the neon fill (Sec.~\ref{ea:sec:ei}). There the collisions are classical, and the ions'
term rests on an estimate, the larger of Scheuer's classical Gaunt factor and one. The
theory of that term is tested elsewhere. The quantum theory gives the Coulomb logarithm
with no cutoff put in by hand (Sec.~\ref{pd:sec:summary}, Eq.~\eqref{env:eq:Lambda1}),
and with no fitted parameter the ladder of Chapter~\ref{lad:chapter} gives
$0.96\pm0.04$ of the measured absorption of ten nitrogen--hydrogen laser plasmas of
Turnbull et al.\ \cite{p1:Turnbull2023,p1:Turnbull2024} and $1.19\pm0.06$ of four
hydrogen ones (Eq.~\eqref{lad:eq:verdict}). Those plasmas, at temperatures such as the
0.36~keV of the hydrogen point of Sec.~\ref{lad:sec:idea}, are far from the afterglow.

{\noindent}(5) \emph{The solvers.} Two independently written solvers of the two-term
equation agree within 0.2\,\% in the electrons' mean energy and in the absorbed power
(Sec.~\ref{qib:sec:solvers}), and the kinetic solver of Chapter~\ref{nl:chapter}
reproduces Langdon's absorption fit within 1.8\,\% (Sec.~\ref{nl:sec:checks}).

\medskip
\noindent\textbf{Standing: verified where the atoms absorb.} The kernel holds to 1.6\,\%
in the THz band, and krypton's cross section is tested to a few per cent for electrons
above 0.95~eV and to about 20\,\% at the afterglow's 0.63~eV. In a thermalized
afterglow above $10^{20}$~m$^{-3}$ the ions take over, and their term rests there on a
classical estimate.

%================================================================
\section{L3: heating}
\label{p1:sec:L3}
%================================================================

{\noindent}(1) In a laser gas the field takes over the electron temperature once its
heating exceeds the electrons' collisional losses, at the intensity $I_{\rm heat}$ of
Eq.~\eqref{nl:eq:thresholds}, below which a wave crosses a weakly ionized plasma
without changing it appreciably \cite{p1:GinzburgGurevich1960}. In the Ne fill's
afterglow the electrons lose $3.6\times10^{8}$~eV\,s$^{-1}$ each at 0.63~eV, and
$I_{\rm heat}$ is $2.8\times10^{2}$~W\,cm$^{-2}$ at 2.45~GHz,
$1.4\times10^{3}$~W\,cm$^{-2}$ at 170~GHz and $1.5\times10^{4}$~W\,cm$^{-2}$ at
0.67~THz (Table~\ref{nl:tab:regime}). Gyrotrons deliver 1~MW at 170~GHz
\cite{p1:Sakamoto2007} and 0.21~MW at 0.67~THz \cite{p1:Glyavin2012}
(Fig.~\ref{nl:fig:regime}), and the second holds $1.5\times10^{4}$~W\,cm$^{-2}$ over
14~cm$^2$.

{\noindent}(2) The power per volume is small against the beam's. To hold the
afterglow's electrons at 0.63~eV the field must give each $3.6\times10^{8}$~eV\,s$^{-1}$,
5.8~kW\,cm$^{-3}$ at $10^{20}$~m$^{-3}$, and to hold them at the beam-on 1.5~eV
$1.1\times10^{9}$~eV\,s$^{-1}$ (Table~\ref{nl:tab:regime}), 18~kW\,cm$^{-3}$. These are
1.0 and 2.9\,\% of the 0.6~MW\,cm$^{-3}$ of beam that makes the beam-on density in a
full kinetic model of the laser (Sec.~\ref{nl:sec:share}).

{\noindent}(3) The heating stays in the regime the kinetic equation describes.
Langdon's flattening of the distribution needs 33 to 51 times $I_{\rm heat}$ in the
afterglow, and Silin's regime a further 11.5 times (Sec.~\ref{nl:sec:regime}). At every
point of the scans the quiver velocity is at most 0.13 of the thermal one
(Sec.~\ref{nl:sec:kinetic}). For electrons whose energy far exceeds the photon's, the
exchange of several photons changes how the energy is gained but not how fast, so the
classical kernel holds at any intensity, and with elastic losses alone a 0.5~eV rise of
the electron temperature in pure krypton needs 53~W\,cm$^{-2}$ at 0.3~THz and
585~W\,cm$^{-2}$ at 1~THz (Sec.~\ref{ea:sec:strong}).

{\noindent}(4) The density answers to the heating. Heated hard, the afterglow ionizes
itself. At $10^{14}$~cm$^{-3}$ and 300~kW\,cm$^{-3}$ the density e-folds in 2.4 to
5.2~ns in pure krypton and in krypton with F$_2$ at 0.1~atm
(Sec.~\ref{env:sec:particles}), and the Ne fill multiplies its electrons more than
$e$-fold in 1~$\mu$s at every pressure at that density (Table~\ref{qib:tab:fbest}).
Above the ionization balance the reheated afterglow becomes a discharge sustained by
the field. A self-sustained discharge forms 0.337 to 0.354 KrF$^*$ per absorbed joule
(Sec.~\ref{nl:sec:wallplug}), and the notes do not address its stability. Below that
balance the F$_2$ removes the electrons. In krypton with 2\,\% F$_2$ attachment takes
them in 0.1 to 17~ns (Sec.~\ref{env:sec:f2}), so reheating acts on electrons that turn
over in nanoseconds.

\medskip
\noindent\textbf{Standing: verified.} $1.4\times10^{3}$ to
$1.5\times10^{4}$~W\,cm$^{-2}$ hold the afterglow's electrons at their temperature, and
a few per cent of the beam's power per volume holds them at the beam-on temperature, in
a regime that the two-term kinetic equation and the classical kernel describe.

%================================================================
\section{What is invalidated}
\label{p1:sec:invalid}
%================================================================

{\noindent}(1) \emph{The laser's own field.} At 10~MW\,cm$^{-2}$, a few saturation
intensities, the KrF field is $7.1\times10^{3}$ times below $I_{\rm heat}$ in the Ne
fill and $1.5\times10^{4}$ times in the Ar fill (Sec.~\ref{nl:sec:regime}). The laser's
light does not reheat its own electrons.

{\noindent}(2) \emph{The nonlinearity of IB.} An auxiliary field reaches only the
nonlinearity of a weakly ionized gas, in which the field sets the electron temperature,
long before Langdon's flattening \cite{p1:Langdon1980} or Silin's saturation, which need
33 to 80 and 380 to 920 times more intensity (Sec.~\ref{nl:sec:answer}). The
nonlinearity of IB is not by itself a lever on the efficiency of a gas laser
(Sec.~\ref{nl:sec:verdict}).

{\noindent}(3) \emph{The frequency as a lever, with the beam on.} A field far above the
collision frequency heats each electron in proportion to its collision frequency and
gives a distribution whose tail reaches krypton's 9.9~eV levels, where a DC field heats
the slow electrons most (Sec.~\ref{nl:sec:share}). At the beam-on density the
electron--electron collisions supply most of that tail anyway, and what is left of the
frequency's advantage is 2 to 7\,\% where the field's power exceeds the beam's
(Sec.~\ref{nl:sec:verdict}). In the afterglow, with 27 times fewer electrons, the lever
survives, which L4 takes up.

%================================================================
\section{L4: partition}
\label{p1:sec:L4}
%================================================================

{\noindent}(1) The formation efficiency $\eta_{\rm form}$ of Eq.~\eqref{nl:eq:etaform}
counts the KrF$^*$ that a partition of the absorbed power forms per absorbed energy, at
the photon energy of 4.997~eV, with the branchings of a KrF kinetic network: 0.87 for
Kr$^*$ through the harpoon in the Ne fill, 0.81 in the Ar fill, and 0.68 through the
ions. The electron beam forms 0.198 in the Ne fill by the same count once its
sub-excitation electrons are credited, 0.175 without them, and a full kinetic model of
the laser gives 0.194 (Sec.~\ref{nl:sec:wallplug}).

{\noindent}(2) At the afterglow's density the frequency still matters. At
$10^{20}$~m$^{-3}$ a field far above the collision frequency, 28.3~THz in the scans,
gives krypton 66.4\,\% of the absorbed power in the Ne fill at the lowest power, against
42.7\,\% from a DC field of the same absorbed power, a ratio of 1.556, and the ratio is
1.596 in the Ar fill (Table~\ref{nl:tab:share}). With electrons sustained by the beam
it is 1.709 and 1.899 (Table~\ref{nl:tab:source}). In the afterglow a 2.45~GHz wave is a
DC field for this purpose, $\langle\nu_m\rangle/\omega=32$, while a THz wave lies above
the collision frequency, $\langle\nu_m\rangle/\omega=0.12$ already at 0.67~THz
(Table~\ref{nl:tab:regime}).

{\noindent}(3) The high frequency forms 0.274 to 0.352 KrF$^*$ per absorbed joule in the
Ne fill at $10^{20}$~m$^{-3}$: 0.274 and 0.351 at the powers of 2 and 5~Td with the growth
model (Table~\ref{nl:tab:share}), and 0.298 and 0.352 for the field's marginal power with
beam-sustained electrons (Table~\ref{nl:tab:source}). That is 1.4 to 1.8 times the
beam's 0.198. A DC field forms 0.175 to 0.326 there. In the Ar fill the high frequency
forms 0.342 to 0.375 against 0.101 to 0.208.

{\noindent}(4) The count leaves out what a reheated afterglow does most.
\begin{itemize}
\item \emph{The electron collisions of Kr$^*$ and KrF$^*$.} At the beam-on density the
  stepwise ionization $e+{\rm Kr}^*\to{\rm Kr}^++2e$ runs at 0.39 to 1.63 of the harpoon
  rate in the Ne fill, at the rate of Kannari et al.\ \cite{p1:Kannari1985}, and lowers
  the KrF$^*$ yield of a Kr$^*$ by 4.8 to 11.8\,\%, which lowers field pumping, mostly
  through Kr$^*$, more than beam pumping, mostly through ions. In a full kinetic model of
  microwave co-pumping the sign of the laser energy that the field gains rests on that
  one rate, and the electron quench of KrF$^*$ is a candidate second
  (Sec.~\ref{nl:sec:wallplug}). No measurement of the stepwise ionization was found
  (Sec.~\ref{qib:sec:ladder}). In the afterglow these collisions weigh more, since
  Kr$^*$ accumulates once its harpoon partner is gone, 10 and 8 times in the pulses of
  Sec.~\ref{qib:sec:pulse}.
\item \emph{The 5s to 5p cycle.} The electrons lift Kr$^*$ to 5p with the measured cross
  sections of Jung et al.\ \cite{p1:Jung2005}, and in the krypton fill at 1~atm and
  $10^{15}$~cm$^{-3}$, at 10~THz, the 5p manifold returns 84\,\% by collisions with the
  gas, 15\,\% superelastically and 1.2\,\% by radiation (Sec.~\ref{qib:sec:ladder}). Net
  of what they give back, the excited atoms take 0.007 to 2.9 times the absorbed power
  from the electrons in the Ne fill and up to 23.5 times in the krypton fill
  (Sec.~\ref{qib:sec:results}). A faster collisional return ends more of the cycle's
  energy in the gas (Sec.~\ref{qib:sec:limits}), and one measured total loss rate of a
  5p level bounds it \cite{p1:Gazeli2021}.
\item \emph{The F$_2$.} Kr$^*$ forms KrF$^*$ only while F$_2$ lasts, and the electrons
  burn it by attachment, $3\times10^{8}$ to $6\times10^{8}$ molecules per electron per
  second in the pulses of Sec.~\ref{qib:sec:pulse}, where half of it is gone at 134~ns at
  $10^{15}$~cm$^{-3}$ and at 11~ns at $10^{16}$~cm$^{-3}$. The afterglow inherits what
  the pump leaves.
\end{itemize}

\medskip
\noindent\textbf{Standing: open.} Per absorbed joule a THz wave in the afterglow would
form 1.4 to 1.8 times the KrF$^*$ of the beam, and a microwave less. The count leaves
out the stepwise ionization, the electron quench of KrF$^*$ and the 5s to 5p cycle,
which grow with the electron temperature and with Kr$^*$, and whether Kr$^*$ becomes
KrF$^*$ at all depends on the F$_2$ the pump leaves.

%================================================================
\section{L5: the laser and the wall plug}
\label{p1:sec:L5}
%================================================================

{\noindent}(1) KrF$^*$ formed in the afterglow is laser energy only if photons extract
it before it is lost. No note follows the photons. The excimers' own kinetics, the
photons, the gas temperature, the atomic fluorine and the walls are left to a
time-domain model of the laser (Sec.~\ref{qib:sec:stands}).

{\noindent}(2) Field pumping beats beam pumping at the wall plug when
$\eta_{\rm src}\eta_{{\rm form},f}>\eta_{\rm beam}\eta_{{\rm form},b}$, that is when the
source's efficiency exceeds 0.57 to 0.65 of the efficiency with which the pulsed power
puts energy into the gas through the foil, for the field's power at the beam-on density
(Sec.~\ref{nl:sec:wallplug}). In the afterglow the KrF$^*$ must also be extracted. With
$\eta_x$ the fraction of the KrF$^*$ that leaves as laser light, the condition is
\begin{equation}
\eta_{\rm src}\,\eta_{{\rm form},f}\,\eta_{x,f}>\eta_{\rm beam}\,\eta_{{\rm form},b}\,\eta_{x,b},
\label{p1:eq:wallplug}
\end{equation}
with $f$ for the field and $b$ for the beam. With the afterglow's
$\eta_{{\rm form},f}=0.274$ to 0.352 against 0.198, the source must reach 0.56 to 0.72
of $\eta_{\rm beam}$ times $\eta_{x,b}/\eta_{x,f}$. No note computes either $\eta_x$.

{\noindent}(3) The sources that the notes cite are the 1~MW gyrotron at 170~GHz and the
0.21~MW gyrotron at 0.67~THz. In the afterglow the first heats 2.5~mm and the second
2.6~cm (L1). The notes cite no source of comparable power at the 1.3~THz that heats
10~cm of the Ne fill's afterglow.

\medskip
\noindent\textbf{Standing: open.} Whether reheating raises the laser's efficiency
depends on the extraction of the afterglow's KrF$^*$ and on the efficiency of a THz
source, neither of which Part~\ref{p1:part:notes} computes.


\chapter{What would settle it}
\label{p1:chap:settle}

{\noindent}(1) The open links hinge on two rates, one return rate and one calculation.
In the order of their cost:
\begin{enumerate}
\item \emph{The stepwise ionization $e+{\rm Kr}^*\to{\rm Kr}^++2e$ and the electron
  quench of KrF$^*$.} They decide the sign of the field's gain (L4). A THz pulse sets the
  electron temperature on demand and its transmission reads the density and the
  temperature back (Secs.~\ref{nl:sec:answer} and \ref{nl:sec:verdict}). With elastic
  losses alone a 0.5~eV rise in krypton needs 53~W\,cm$^{-2}$ at 0.3~THz
  (Sec.~\ref{ea:sec:strong}).
\item \emph{The collisional return of 5p to 5s in krypton, level by level}
  \cite{p1:Whitehead1995}. It is the largest uncertainty of the best frequency at the
  dense krypton points, 13\,\% in it and 19\,\% in the electron temperature
  (Sec.~\ref{qib:sec:limits}), and it decides how much of the 5s to 5p cycle's energy
  ends in the gas (L4).
\item \emph{Drift velocities in krypton below 0.95~eV}, across the Ramsauer minimum where
  the afterglow's electrons sit and which one experiment covers
  (Sec.~\ref{env:sec:sources}). They would test the absorption at 0.63~eV to the few per
  cent of the warmer electrons (L2).
\item \emph{A time-dependent kinetic model of the laser with the THz absorption computed
  from its running state}, the electron collisions of the excited atoms of
  Table~\ref{qib:tab:erows}, the burn of the F$_2$ and the extraction, run in pairs with
  and without the field and with the uncertain rates bracketed
  (Secs.~\ref{qib:sec:stands} and \ref{nl:sec:verdict}). It gives $\eta_x$ and the laser
  energy per absorbed joule, the two numbers that L5 lacks.
\item \emph{The wall plug}, Eq.~\eqref{p1:eq:wallplug}, with those numbers and the
  efficiency of a source near 1.3~THz.
\end{enumerate}

{\noindent}(2) The proposal is invalidated if the run of item~4 finds the laser energy
that the field gains negative over the measured range of the two rates of item~1, if
the pump leaves too little F$_2$ for the reheated Kr$^*$ to form KrF$^*$, or if no source
meets Eq.~\eqref{p1:eq:wallplug} with the extraction that the run finds.

\begin{thebibliography}{99}
\raggedright
\bibitem{p1:Nakamura1991} Y.~Nakamura, Meeting Abstracts of the Physical Society of Japan,
1991 autumn meeting, abstract 27p-ZA-3, p.~103 (1991),
doi:\href{https://doi.org/10.11316/jpsgaiyoa.1991.4.0_103}{10.11316/jpsgaiyoa.1991.4.0\_103}.
\bibitem{p1:Pack1962} J.~L.~Pack, R.~E.~Voshall and A.~V.~Phelps, Phys. Rev. \textbf{127},
2084 (1962), doi:\href{https://doi.org/10.1103/PhysRev.127.2084}{10.1103/PhysRev.127.2084}.
\bibitem{p1:Turnbull2023} D.~Turnbull \emph{et al.}, Phys. Rev. Lett. \textbf{130}, 145103
(2023), doi:\href{https://doi.org/10.1103/PhysRevLett.130.145103}{10.1103/PhysRevLett.130.145103}.
\bibitem{p1:Turnbull2024} D.~Turnbull \emph{et al.}, Phys. Plasmas \textbf{31}, 063304 (2024),
doi:\href{https://doi.org/10.1063/5.0203446}{10.1063/5.0203446}.
\bibitem{p1:GinzburgGurevich1960} V.~L.~Ginzburg and A.~V.~Gurevich, Sov. Phys. Usp.
\textbf{3}, 115 and 175 (1960),
doi:\href{https://doi.org/10.1070/PU1960v003n01ABEH003261}{10.1070/PU1960v003n01ABEH003261} and
doi:\href{https://doi.org/10.1070/PU1960v003n02ABEH003267}{10.1070/PU1960v003n02ABEH003267}.
\bibitem{p1:Sakamoto2007} K.~Sakamoto \emph{et al.}, Nature Phys. \textbf{3}, 411 (2007),
doi:\href{https://doi.org/10.1038/nphys599}{10.1038/nphys599}.
\bibitem{p1:Glyavin2012} M.~Yu.~Glyavin, A.~G.~Luchinin, G.~S.~Nusinovich, J.~Rodgers,
D.~G.~Kashyn, C.~A.~Romero-Talamas and R.~Pu, Appl. Phys. Lett. \textbf{101}, 153503 (2012),
doi:\href{https://doi.org/10.1063/1.4757290}{10.1063/1.4757290}.
\bibitem{p1:Langdon1980} A.~B.~Langdon, Phys. Rev. Lett. \textbf{44}, 575 (1980),
doi:\href{https://doi.org/10.1103/PhysRevLett.44.575}{10.1103/PhysRevLett.44.575}.
\bibitem{p1:Kannari1985} F.~Kannari, M.~Obara and T.~Fujioka, J. Appl. Phys. \textbf{57},
4309 (1985), doi:\href{https://doi.org/10.1063/1.334590}{10.1063/1.334590}.
\bibitem{p1:Jung2005} R.~O.~Jung, T.~E.~Stone, J.~B.~Boffard, L.~W.~Anderson and C.~C.~Lin,
Phys. Rev. Lett. \textbf{94}, 163202 (2005),
doi:\href{https://doi.org/10.1103/PhysRevLett.94.163202}{10.1103/PhysRevLett.94.163202}.
\bibitem{p1:Gazeli2021} K.~Gazeli, X.~Aubert, S.~Prasanna, C.~Y.~Duluard, G.~Lombardi and
K.~Hassouni, Phys. Plasmas \textbf{28}, 043301 (2021),
doi:\href{https://doi.org/10.1063/5.0041471}{10.1063/5.0041471}.
\bibitem{p1:Whitehead1995} C.~A.~Whitehead, H.~Pournasr, M.~R.~Bruce, H.~Cai, J.~Kohel,
W.~B.~Layne and J.~W.~Keto, J. Chem. Phys. \textbf{102}, 1965 (1995),
doi:\href{https://doi.org/10.1063/1.468763}{10.1063/1.468763}.
\end{thebibliography}
